\section{Future directions}
\label{sec:future-directions}

\begin{question}[Non-simply-connected detection]
\label{conj:non-simply-connected-detection}
Do there exist a semisimple premodular fusion category $\cC$ over $\mathbb C$,
a monoidal natural automorphism $\eta$ of its identity functor, homeomorphic
but non-diffeomorphic closed oriented non-simply-connected smooth
4-manifolds $X$ and $Y$, an orientation-preserving homeomorphism
$f\colon X\to Y$, and a class $\alpha\in H_2(X;\mathbb Z)$ for which
\[
  \CY_{\cC,\eta}\bigl([(X,\alpha)]\bigr)
  \ne\CY_{\cC,\eta}\bigl([(Y,f_*\alpha)]\bigr).
\]
\end{question}

Gompf proved that any two smoothings of a compact orientable 4-manifold
become diffeomorphic after connected sum with copies of $S^2\times S^2$
\cite{Gompf1984}.  That unmarked theorem does not by itself identify the
induced map on $H_2$.  The precise missing input for our argument is the
following marked version: given $f\colon X\to Y$ as above, does there exist
$r\geq0$ and an orientation-preserving diffeomorphism
\[
 F\colon X\#r(S^2\times S^2)\longrightarrow
          Y\#r(S^2\times S^2)
\]
such that, under the canonical connected-sum decompositions,
\[
 F_*=f_*\oplus\operatorname{id}_{rH}
 \quad\text{on }H_2?
\]
An affirmative answer would extend the stabilization-and-cancellation proof
of \cref{cor:no-simply-connected-exotic-detection} to this setting and give a
negative answer to the preceding question.  We make neither assertion here.

\begin{conjecture}[Full coefficient-group extension]
\label{conj:full-coefficient-extension}
Let \(\cC\) be a premodular fusion category and let
\[
  A_{\cC}=\operatorname{Aut}_{\otimes}(\operatorname{Id}_{\cC}).
\]
For every closed, connected, oriented smooth 4-manifold \(X\), the
fixed-character construction should extend to an invariant
\[
  \CY_{\cC}(X,-)\colon H_2(X;A_{\cC})\longrightarrow\mathbb C.
\]
For a fixed character $\eta\in A_{\cC}$, let
\[
  \varphi_\eta\colon\mathbb Z\longrightarrow A_{\cC},
  \qquad \varphi_\eta(n)=\eta^n.
\]
The pullback of $\CY_{\cC}(X,-)$ along the induced coefficient map
\[
  (\varphi_\eta)_*\colon H_2(X;\mathbb Z)\longrightarrow H_2(X;A_{\cC})
\]
should be the function
$\alpha\mapsto\CY_{\cC,\eta}\bigl([(X,\alpha)]\bigr)$ constructed in this
paper.
\end{conjecture}

The coefficient group $A_{\cC}$ and its action on simple objects are those
of \cref{prop:characters-universal-grading}; this is also the coefficient
system used in \cite[Definitions~1 and~9]{Yetter1993}.  Establishing the
extension requires redoing representative independence and every
character-aware attach--transport--detach calculation with arbitrary
$A_{\cC}$-labels rather than powers of one fixed character.  In particular,
the resulting numerical formulas would require the family of twisted Gauss
sums
\[
  \tau^{\pm}(a)
  :=\sum_{i\in\operatorname{Irr}(\cC)}
       d_i^2\theta_i^{\pm1}\chi_i(a),
  \qquad a\in A_{\cC},
\]
where $\chi_i(a)$ is the scalar by which $a$ acts on $i$.  Neither this full
coefficient-group invariant nor its simply connected closed formulas are
claimed here.
